\section{Introduction}
\label{sec:introduction}

The boundary of a useful model no longer coincides with one checkpoint. A
deployment may keep latency-sensitive or private work on a robot or edge NPU,
draw on specialists in a sovereign data center, and reach a frontier cloud
service only when the task warrants it. Open and closed models, accelerator
generations, network conditions, prices, and jurisdictions evolve on different
schedules. To the caller, however, the system should still behave as one model:
one interface, one set of constraints, and one attributable result.

Heterogeneity is also an opportunity for conditional capacity. When constituent
errors or resource profiles are complementary, a bounded controller can
outperform a fixed choice under a declared objective. A cost-first variant
avoids unnecessary work; an accuracy-first variant spends a declared budget on
disagreement detection, verification, and synthesis. Security-first and
grounding-first variants change which traces are admissible, not merely how
they are ranked. Mixture-of-Experts (MoE) established conditional computation
within a checkpoint by sparsely activating jointly trained experts
\citep{shazeer_2017_moe}. We move the boundary outward. A \emph{Mixture-of-Models}
(MoM) is a versioned composite model whose lifecycle engine realizes requests
through preference-conditioned, resource-bounded traces over independently
versioned constituents. Routing or selection is one possible topology, not the
architecture itself. MoM is therefore a form of system-level intelligence:
capability is not localized in any constituent, but emerges from a control
policy that allocates heterogeneous model and system resources under workload,
preference, and deployment state. This view extends the Workload--Router--Pool
thesis that workload, control, and physical realization are coupled design
variables rather than independent layers \citep{chen_2026_wrp}.

Routers, cascades, ensembles, and agent pipelines already exploit parts of this
space \citep{chen_2023_frugalgpt,ong_2024_routellm,wang_2024_moa,
kandogan_2024_compound}. In these systems, the composite is typically treated
and evaluated as a routing, ensemble, or application pipeline rather than
versioned and evaluated as one model artifact. For sovereign AI deployments,
this boundary is consequential:
concrete cloud-sovereignty requirements can make a capable call inadmissible
because of residency, audit, or provider control
\citep{european_commission_2026_sovereignty}. We therefore place MoM at the
model-invocation infrastructure layer, below agent harnesses. Agents retain
application goals and tool intent; a complete lifecycle engine owns
admissibility, finite composition, budgets, and attribution. \vsr{} supplies an
initial execution substrate at this boundary. We use it to ground the
architecture while proposing the portable identity, binding, and conformance
layer needed to build, train, evaluate, exchange, and serve composite models as
durable model objects.

\begin{figure}[H]
  \centering
  \resizebox{\linewidth}{!}{\begin{tikzpicture}[
  x=1cm,y=1cm,font=\sffamily,>=Latex,
  flow/.style={-{Latex[length=2.0mm]},draw=MoMGray,line width=0.65pt},
  bind/.style={-{Latex[length=1.7mm]},draw=MoMGray,densely dashed,line width=0.55pt},
  boundary/.style={draw=MoMGray,fill=MoMBlueLight!9,densely dashed,
    rounded corners=3pt,line width=0.6pt},
  card/.style={draw=MoMBorder,fill=white,rounded corners=2pt,line width=0.55pt,
    minimum height=0.68cm,inner xsep=5pt,align=center,
    font=\sffamily\fontsize{7.2}{8.1}\selectfont},
  soft/.style={card,draw=none,fill=MoMBlueLight},
  core/.style={card,draw=none,fill=MoMDark,text=white,
    minimum height=0.74cm,font=\sffamily\bfseries\fontsize{7.2}{8.1}\selectfont},
  active/.style={card,draw=MoMDark,fill=MoMBlue,line width=0.75pt},
  tiny/.style={font=\sffamily\fontsize{6.7}{7.5}\selectfont,text=MoMMuted},
  grouptitle/.style={font=\sffamily\bfseries\fontsize{7.1}{8.0}\selectfont,
    text=MoMText},
  paneltitle/.style={font=\sffamily\bfseries\fontsize{8.4}{9.4}\selectfont,
    text=MoMText}
]

% Symmetric panels compare the two loci of conditional computation.
\draw[draw=MoMBorder,rounded corners=4pt,line width=0.65pt]
  (0,0.15) rectangle (8.0,5.60);
\draw[draw=MoMBorder,rounded corners=4pt,line width=0.65pt]
  (8.30,0.15) rectangle (16.30,5.60);
\node[paneltitle,anchor=west] at (0.32,5.28) {(a) Mixture-of-Experts};
\node[paneltitle,anchor=west] at (8.62,5.28) {(b) Mixture-of-Models};

% MoE: gate and experts share one training and checkpoint boundary.
\draw[boundary] (1.45,1.28) rectangle (6.25,4.02);
\node[grouptitle] at (3.85,4.43) {one jointly trained checkpoint};

\node[soft,minimum width=1.12cm] (token) at (0.72,2.72) {token};
\node[core,minimum width=1.22cm] (gate) at (2.25,2.72) {gate};
\node[active,minimum width=1.78cm] (e1) at (4.55,3.38) {expert 1};
\node[card,minimum width=1.78cm] (e2) at (4.55,2.72) {expert 2};
\node[card,minimum width=1.78cm] (e3) at (4.55,2.06) {expert 3};
\node[soft,minimum width=1.02cm] (state) at (7.15,2.72) {state};

\draw[flow] (token) -- (gate);
\draw[flow] (gate.east) -- ++(0.34,0) |- (e1.west);
\draw[draw=MoMBorder,line width=0.45pt] (gate.east) -- ++(0.34,0) |- (e2.west);
\draw[draw=MoMBorder,line width=0.45pt] (gate.east) -- ++(0.34,0) |- (e3.west);
\draw[flow] (e1.east) -- ++(0.32,0) |- (state.west);
\draw[draw=MoMBorder,line width=0.45pt] (e2.east) -- ++(0.32,0) |- (state.west);
\draw[draw=MoMBorder,line width=0.45pt] (e3.east) -- ++(0.32,0) |- (state.west);
\node[tiny] at (3.85,0.72) {token-level sparse activation};

% MoM: engine and constituent models form the complete virtual model.
\draw[boundary] (9.65,1.28) rectangle (15.02,4.02);
\node[grouptitle] at (12.34,4.43) {one virtual composite model};

\node[soft,minimum width=1.12cm] (request) at (8.98,3.30) {request};
\node[card,minimum width=1.22cm] (profile) at (8.98,2.12)
  {preferences\\[-1pt]constraints};
\node[core,minimum width=1.34cm] (engine) at (10.55,2.72) {Engine};
\node[active,minimum width=1.82cm] (m1) at (13.05,3.38) {local model};
\node[card,minimum width=1.82cm] (m2) at (13.05,2.72) {private model};
\node[active,minimum width=1.82cm] (m3) at (13.05,2.06) {cloud model};
\node[soft,minimum width=1.08cm] (answer) at (15.74,2.72) {answer};

\draw[flow] (request) -- (engine);
\draw[flow] (profile) -- (engine);
\draw[flow] (engine.east) -- ++(0.34,0) |- (m1.west);
\draw[draw=MoMBorder,line width=0.45pt] (engine.east) -- ++(0.34,0) |- (m2.west);
\draw[flow] (engine.east) -- ++(0.34,0) |- (m3.west);
\draw[flow] (m1.east) -- ++(0.32,0) |- (answer.west);
\draw[draw=MoMBorder,line width=0.45pt] (m2.east) -- ++(0.32,0) |- (answer.west);
\draw[flow] (m3.east) -- ++(0.32,0) |- (answer.west);

% Physical resources remain outside the logical composite-model boundary.
% A single realization card keeps the logical-to-physical edge unambiguous.
\node[card,minimum width=2.35cm,minimum height=0.58cm] (fleet) at (13.05,0.57)
  {{\bfseries physical fleet}\\[-1pt]
   {\fontsize{6.3}{7.0}\selectfont\color{MoMMuted} GPU / CPU / NPU}};
\draw[bind] (13.05,1.28) -- node[tiny,anchor=west,xshift=2pt] {bind} (fleet.north);

\end{tikzpicture}
}
  \caption{\textbf{Conditional computation at two model boundaries.} In MoE,
  a gate selects experts inside one jointly trained checkpoint. In MoM, an
  engine selects and composes independently versioned models behind one model
  interface; a deployment binding maps their logical identities onto a physical
  fleet. Highlighted nodes show one realized trace.}
  \label{fig:moe-vs-mom}
\end{figure}

Crossing the model boundary changes the optimization object. MoM does not
require a shared parameter space. Instead, it learns or searches over
admissible traces, then promotes changes to control state, pools, topology,
contracts, or deployment as explicit revisions. Candidates are evaluated on
realized utility and constraint outcomes. Uniform utilization is not an
intrinsic goal because constituents differ in capability, price, capacity, and
policy eligibility.

This paper makes three contributions:
\begin{itemize}
  \item We formalize MoM as a preference-conditioned distribution over finite,
  typed execution traces, separating hard admissibility from soft objectives
  and defining capacity diagnostics and failure criteria.
  \item We specify a composite-model format and lifecycle from logical
  specification to portable bundle, realization profile, environment binding,
  resolution lock, and attributable execution.
  \item We ground the execution and systems claims with released multi-model
  runs and a routing-aware capacity case study, and state the matched-budget
  tests required to establish conditional scaling.
\end{itemize}
